\section*{Stavová tabulka k právu a strategiím}
\addcontentsline{toc}{section}{Stavová tabulka k právu a strategiím}
\begin{longtable}{p{0.25\linewidth}p{0.24\linewidth}p{0.20\linewidth}p{0.24\linewidth}}
\toprule
\textbf{Předpis / dokument} & \textbf{Stav k 18. 6. 2026} & \textbf{Použitelnost / účinnost} & \textbf{Co si pamatovat}\\
\midrule
NIS2 - směrnice 2022/2555 & Platná; transpoziční lhůta skončila 17. 10. 2024. & V ČR transponována hlavně zákonem o kybernetické bezpečnosti (č. 264/2025 Sb.). & Směrnice se neuplatní jako běžný český zákon sama o sobě; do praxe ji převádí národní úprava.\\
\addlinespace
Zákon o kybernetické bezpečnosti (č. 264/2025 Sb.) & Účinný, nahradil starý zákon o kybernetické bezpečnosti z roku 2014. & 1. 11. 2025. & Režim vyšších a nižších povinností, sebeidentifikace regulované služby, hlášení incidentů, odpovědnost vedení.\\
\addlinespace
Vyhlášky k novému ZKB & Účinné od 1. 11. 2025. & Zejména vyhláška o regulovaných službách, vyhláška o bezpečnostních opatřeních pro vyšší režim a vyhláška o bezpečnostních opatřeních pro nižší režim. & U zkoušky je důležité rozlišit, jestli subjekt vůbec poskytuje regulovanou službu a v jakém režimu.\\
\addlinespace
Národní strategie kybernetické bezpečnosti 2026-2030 & Zveřejněna NÚKIB 8. 9. 2025. & Strategický dokument od roku 2026. & Navazuje na Bezpečnostní strategii ČR 2023, Obrannou strategii 2023, EU a NATO dokumenty; akční plán je samostatný implementační dokument.\\
\addlinespace
CER - směrnice 2022/2557 a zákon o odolnosti subjektů kritické infrastruktury (č. 266/2025 Sb.) & Nová úprava odolnosti kritických subjektů. & Zákon účinný od 19. 8. 2025; k 18. 6. 2026 nebylo v dostupných oficiálních zdrojích potvrzeno účinné nařízení vlády o základních službách a kritériích významnosti. & Vedle kybernetické bezpečnosti řeší fyzickou a organizační odolnost; plné nové určování subjektů závisí na prováděcím nařízení a rozhodnutích MV.\\
\addlinespace
DORA - nařízení 2022/2554 & Platné a použitelné. & 17. 1. 2025. & Lex specialis pro finanční sektor: ICT rizika, incidenty, testování digitální odolnosti, dohled nad ICT třetími stranami.\\
\addlinespace
CRA - nařízení 2024/2847 & V platnosti. & Hlášení od 11. 9. 2026, hlavní povinnosti od 11. 12. 2027. & Produkty s digitálními prvky: bezpečný vývoj, zranitelnosti, aktualizace, CE logika.\\
\addlinespace
EU e-evidence - nařízení 2023/1543 & Platné. & Použitelné od 18. 8. 2026. & EPOC je evropský příkaz k vydání, EPOC-PR je příkaz k uchování. Neplést lhůty vydání a uchování.\\
\addlinespace
EU e-evidence - směrnice 2023/1544 & Platná; členské státy měly transponovat do 18. 2. 2026. & Navazuje na použitelnost e-evidence režimu. & Určené provozovny a právní zástupci poskytovatelů pro příjem příkazů.\\
\addlinespace
eIDAS 2.0 - nařízení 2024/1183 & V platnosti od 20. 5. 2024. & EUDI peněženky se odvíjejí od prováděcích aktů; obecně se uvádí horizont roku 2026, ale přesné lhůty je třeba ověřit podle aktuálních aktů. & Kvalifikovaný elektronický podpis má účinky vlastnoručního podpisu; EUDI peněženka přidává digitální identitu.\\
\addlinespace
2. dodatkový protokol k Budapešťské úmluvě - CETS 224 & Stále není v platnosti, protože je potřeba pět ratifikací. & Vstup po splnění ratifikační podmínky. & Nelze ho uvádět jako už účinný právní základ; je to budoucí nástroj přímější spolupráce.\\
\addlinespace
Úmluva OSN proti kyberkriminalitě & Není v platnosti; UNTC uvádí 77 signatářů a 3 strany. & Vstup 90 dní po 40. ratifikaci. & Podpis není ratifikace. U zkoušky říct, že jde zatím hlavně o politický a budoucí kooperační rámec.\\
\bottomrule
\end{longtable}

\newpage

